\documentclass[12pt,letterpaper]{article}

% ------------------------------------------------
% Paquetes
% ------------------------------------------------
\usepackage[utf8]{inputenc}
\usepackage[T1]{fontenc}
\usepackage[spanish]{babel}

\usepackage{amsmath,amssymb,amsthm,mathtools,bm}
\usepackage{booktabs}
\usepackage{array}
\usepackage{geometry}
\usepackage{setspace}
\usepackage{microtype}
\usepackage{parskip}
\usepackage{titlesec}
\usepackage{fancyhdr}
\usepackage{enumitem}
\usepackage{xcolor}
\usepackage{listings}
\usepackage{tcolorbox}
\usepackage{hyperref}

\tcbuselibrary{skins,breakable}

% ------------------------------------------------
% Geometría
% ------------------------------------------------
\geometry{
  top=2.8cm,
  bottom=3cm,
  left=3.2cm,
  right=3.2cm,
  headheight=14pt
}

\setstretch{1.20}

% ------------------------------------------------
% Encabezados
% ------------------------------------------------
\pagestyle{fancy}
\fancyhf{}
\fancyhead[C]{\small\textsc{Econometría --- Gujarati --- Ejercicio 13.8}}
\fancyfoot[C]{\small\thepage}

\renewcommand{\headrulewidth}{0.4pt}

% ------------------------------------------------
% Títulos
% ------------------------------------------------
\titleformat{\section}
{\large\bfseries\scshape}
{\thesection.}{0.6em}{}
[\vspace{-4pt}\rule{\linewidth}{0.4pt}]

\titleformat{\subsection}
{\normalsize\bfseries}
{\thesubsection.}{0.5em}{}

% ------------------------------------------------
% Caja destacada
% ------------------------------------------------
\newtcolorbox{clave}{
  colback=white,
  colframe=black,
  boxrule=0.5pt,
  arc=3pt,
  left=8pt,
  right=8pt,
  top=6pt,
  bottom=6pt,
  breakable
}

% ------------------------------------------------
% Stata style
% ------------------------------------------------
\definecolor{stataback}{RGB}{248,248,248}
\definecolor{stataframe}{RGB}{180,180,180}
\definecolor{statacomment}{RGB}{100,100,100}
\definecolor{statakeyword}{RGB}{0,70,140}

\lstdefinelanguage{Stata}{
  morekeywords={
    clear,set,seed,obs,gen,generate,regress,
    summarize,scalar,display,quietly,
    if,replace,log,close
  },
  sensitive=false,
  morecomment=[l]{*},
  morecomment=[l]{//},
  morestring=[b]"
}

\lstset{
  language=Stata,
  backgroundcolor=\color{stataback},
  basicstyle=\ttfamily\footnotesize,
  keywordstyle=\color{statakeyword}\bfseries,
  commentstyle=\color{statacomment}\itshape,
  frame=single,
  framesep=4pt,
  rulecolor=\color{stataframe},
  xleftmargin=8pt,
  xrightmargin=4pt,
  breaklines=true,
  showstringspaces=false
}

% ------------------------------------------------
% Comandos útiles
% ------------------------------------------------
\newcommand{\Var}{\operatorname{Var}}
\newcommand{\Cov}{\operatorname{Cov}}
\newcommand{\E}{\operatorname{E}}
\newcommand{\plim}{\operatorname{plim}}

% ==========================================================
\begin{document}
% ==========================================================

\begin{center}

{\LARGE\bfseries Ejercicio 13.8 --- Error de Medición}\\[4pt]

{\LARGE\bfseries y Sesgo de Atenuación en MCO}\\[12pt]

{\large\itshape Econometría --- Gujarati \& Porter}\\[10pt]

\rule{0.55\linewidth}{0.6pt}\\[10pt]

{\normalsize
Ejercicio 13.8 resuelto por Emanuel Quintana Silva}\\[4pt]

{\small
Hipótesis de ingreso permanente de Friedman y errores en variables}

\end{center}

\vspace{12pt}

% ==========================================================
\section{Enunciado}
% ==========================================================

Según la hipótesis de ingreso permanente de Friedman:

\begin{equation}
Y_i^{*}=\alpha+\beta X_i^{*}
\end{equation}

donde:

\begin{itemize}
\item $Y_i^{*}$ = gasto de consumo permanente
\item $X_i^{*}$ = ingreso permanente
\end{itemize}

Sin embargo, las variables observadas contienen errores de medición:

\begin{align}
Y_i &= Y_i^{*}+u_i \\
X_i &= X_i^{*}+v_i
\end{align}

donde:

\begin{itemize}
\item $u_i$ = error de medición en $Y_i^{*}$
\item $v_i$ = error de medición en $X_i^{*}$
\end{itemize}

La función observada queda:

\begin{equation}
Y_i=\alpha+\beta(X_i-v_i)+u_i
\end{equation}

o equivalentemente:

\begin{equation}
Y_i=\alpha+\beta X_i+(u_i-\beta v_i)
\end{equation}

Suponiendo:

\begin{enumerate}[leftmargin=1.8em]
\item $\E(u_i)=\E(v_i)=0$
\item $\Var(u_i)=\sigma_u^2$ y $\Var(v_i)=\sigma_v^2$
\item $\Cov(Y_i^{*},u_i)=0$
\item $\Cov(u_i,X_i^{*})
      =
      \Cov(v_i,Y_i^{*})
      =
      \Cov(u_i,v_i)=0$
\end{enumerate}

demostrar que:

\begin{equation}
\plim(\hat{\beta})
=
\frac{\beta}
{1+\left(\frac{\sigma_v^2}{\sigma_X^2}\right)}
\end{equation}

donde:

\begin{equation}
\sigma_X^2=\Var(X_i^{*})
\end{equation}

Además:

\begin{enumerate}[leftmargin=1.8em]
\item[a)] Analizar la naturaleza del sesgo.
\item[b)] Determinar si $\hat{\beta}$ converge al verdadero $\beta$
cuando $n\to\infty$.
\end{enumerate}

% ==========================================================
\section{Modelo observado}
% ==========================================================

El verdadero modelo poblacional es:

\begin{equation}
Y_i^{*}=\alpha+\beta X_i^{*}
\end{equation}

Pero observamos:

\begin{align}
Y_i &= Y_i^{*}+u_i \\
X_i &= X_i^{*}+v_i
\end{align}

Sustituyendo:

\begin{equation}
Y_i
=
\alpha+\beta X_i^{*}+u_i
\end{equation}

y como:

\begin{equation}
X_i^{*}=X_i-v_i
\end{equation}

entonces:

\begin{align}
Y_i
&=
\alpha+\beta(X_i-v_i)+u_i \\[4pt]
&=
\alpha+\beta X_i+(u_i-\beta v_i)
\end{align}

Definimos:

\begin{equation}
w_i=u_i-\beta v_i
\end{equation}

de modo que:

\begin{equation}
Y_i=\alpha+\beta X_i+w_i
\end{equation}

% ==========================================================
\section{Correlación entre regresor y error}
% ==========================================================

El supuesto fundamental de MCO exige:

\begin{equation}
\Cov(X_i,w_i)=0
\end{equation}

Sin embargo:

\begin{align}
X_i &= X_i^{*}+v_i \\
w_i &= u_i-\beta v_i
\end{align}

Por lo tanto:

\begin{align}
\Cov(X_i,w_i)
&=
\Cov(X_i^{*}+v_i,\;u_i-\beta v_i)
\end{align}

Expandiendo:

\begin{align}
&=
\Cov(X_i^{*},u_i)
-\beta\Cov(X_i^{*},v_i)
+\Cov(v_i,u_i)
-\beta\Cov(v_i,v_i)
\end{align}

Usando los supuestos:

\begin{align}
\Cov(X_i^{*},u_i)&=0 \\
\Cov(X_i^{*},v_i)&=0 \\
\Cov(v_i,u_i)&=0
\end{align}

queda:

\begin{equation}
\Cov(X_i,w_i)
=
-\beta\Var(v_i)
\end{equation}

es decir:

\begin{equation}
\boxed{
\Cov(X_i,w_i)
=
-\beta\sigma_v^2
}
\end{equation}

Por lo tanto:

\begin{equation}
\Cov(X_i,w_i)\neq0
\end{equation}

y el estimador MCO es inconsistente.

% ==========================================================
\section{Probabilidad límite del estimador MCO}
% ==========================================================

En muestras grandes:

\begin{equation}
\plim(\hat{\beta})
=
\frac{\Cov(X_i,Y_i)}
{\Var(X_i)}
\end{equation}

Equivalentemente:

\begin{equation}
\plim(\hat{\beta})
=
\beta
+
\frac{\Cov(X_i,w_i)}
{\Var(X_i)}
\end{equation}

Ya obtuvimos:

\begin{equation}
\Cov(X_i,w_i)
=
-\beta\sigma_v^2
\end{equation}

Ahora calculamos $\Var(X_i)$.

Como:

\begin{equation}
X_i=X_i^{*}+v_i
\end{equation}

entonces:

\begin{align}
\Var(X_i)
&=
\Var(X_i^{*}+v_i) \\[4pt]
&=
\Var(X_i^{*})+\Var(v_i)
\end{align}

porque:

\begin{equation}
\Cov(X_i^{*},v_i)=0
\end{equation}

Así:

\begin{equation}
\Var(X_i)
=
\sigma_X^2+\sigma_v^2
\end{equation}

Sustituyendo:

\begin{align}
\plim(\hat{\beta})
&=
\beta
+
\frac{-\beta\sigma_v^2}
{\sigma_X^2+\sigma_v^2}
\end{align}

Factorizando:

\begin{align}
&=
\beta
\left(
1-
\frac{\sigma_v^2}
{\sigma_X^2+\sigma_v^2}
\right)
\end{align}

Simplificando:

\begin{align}
&=
\beta
\left(
\frac{\sigma_X^2}
{\sigma_X^2+\sigma_v^2}
\right)
\end{align}

Finalmente:

\begin{equation}
\boxed{
\plim(\hat{\beta})
=
\frac{\beta}
{1+\left(\frac{\sigma_v^2}{\sigma_X^2}\right)}
}
\end{equation}

que era lo que se quería demostrar.

% ==========================================================
\section{Interpretación económica}
% ==========================================================

El error de medición en la variable explicativa viola el supuesto clásico:

\begin{equation}
\Cov(X_i,u_i)=0
\end{equation}

porque el error de medición $v_i$ aparece tanto:

\begin{itemize}
\item en el regresor observado $X_i$
\item como en el término de perturbación compuesto $w_i$
\end{itemize}

Esto induce correlación entre regresor y error, generando:

\begin{center}
\textbf{sesgo e inconsistencia del estimador MCO.}
\end{center}

% ==========================================================
\section{Inciso a) --- Naturaleza del sesgo}
% ==========================================================

Obtuvimos:

\begin{equation}
\plim(\hat{\beta})
=
\beta
\left(
\frac{\sigma_X^2}
{\sigma_X^2+\sigma_v^2}
\right)
\end{equation}

Como:

\begin{equation}
0<
\frac{\sigma_X^2}
{\sigma_X^2+\sigma_v^2}
<1
\end{equation}

entonces:

\begin{equation}
|\plim(\hat{\beta})|<|\beta|
\end{equation}

Por lo tanto, el estimador está sesgado hacia cero.

\begin{clave}

Este fenómeno se conoce como:

\begin{center}
\textbf{Sesgo de atenuación}
\end{center}

(\textit{attenuation bias})

o también:

\begin{center}
\textbf{Regression dilution bias}
\end{center}

\end{clave}

Si:

\begin{equation}
\beta>0
\end{equation}

entonces:

\begin{equation}
\hat{\beta}<\beta
\end{equation}

Si:

\begin{equation}
\beta<0
\end{equation}

entonces:

\begin{equation}
\hat{\beta}>\beta
\end{equation}

En ambos casos, el efecto estimado se aproxima artificialmente a cero.

% ==========================================================
\section{Inciso b) --- ¿Desaparece el sesgo cuando $n\to\infty$?}
% ==========================================================

No.

Aunque el tamaño muestral aumente indefinidamente:

\begin{equation}
\plim(\hat{\beta})\neq\beta
\end{equation}

porque:

\begin{equation}
\plim(\hat{\beta})
=
\beta
\left(
\frac{\sigma_X^2}
{\sigma_X^2+\sigma_v^2}
\right)
\end{equation}

y este factor es distinto de 1 siempre que:

\begin{equation}
\sigma_v^2>0
\end{equation}

Por consiguiente:

\begin{clave}

El estimador MCO es inconsistente en presencia de errores de medición
en la variable explicativa.

Aumentar el tamaño de muestra no elimina el sesgo.

\end{clave}

Solamente si:

\begin{equation}
\sigma_v^2=0
\end{equation}

(no existe error de medición), entonces:

\begin{equation}
\plim(\hat{\beta})=\beta
\end{equation}

y MCO recupera consistencia.

% ==========================================================
\section{Verificación empírica con Stata}
% ==========================================================

Para ilustrar el sesgo de atenuación simulamos datos donde:

\begin{equation}
Y_i^{*}=1+2X_i^{*}+u_i
\end{equation}

pero observamos:

\begin{equation}
X_i=X_i^{*}+v_i
\end{equation}

La teoría predice:

\begin{equation}
\plim(\hat{\beta})
=
\frac{\beta}
{1+\left(\frac{\sigma_v^2}{\sigma_X^2}\right)}
\end{equation}

% ==========================================================
\subsection{Código Stata}
% ==========================================================

\begin{lstlisting}[caption={Simulación de error de medición y sesgo de atenuación}]
clear
set seed 12345
set obs 5000

* Variable verdadera
gen X_star = rnormal(10,2)

* Error estructural
gen u = rnormal(0,1)

* Error de medicion
gen v = rnormal(0,2)

* Modelo verdadero
gen Y = 1 + 2*X_star + u

* Variable observada con error
gen X = X_star + v

* Regresion usando variable observada
regress Y X

* Regresion usando variable verdadera
regress Y X_star

* Calculo del plim teorico
quietly summarize X_star
scalar sigmaX2 = r(Var)

quietly summarize v
scalar sigmav2 = r(Var)

scalar beta_true = 2

scalar beta_plim = beta_true / ///
    (1 + (sigmav2/sigmaX2))

display _newline ///
"=== Resultado teorico ==="

display ///
"Var(X*)        = " %9.4f sigmaX2

display ///
"Var(v)         = " %9.4f sigmav2

display ///
"Beta verdadero = " %9.4f beta_true

display ///
"plim teorico   = " %9.4f beta_plim
\end{lstlisting}

% ==========================================================
\section{Resultados esperados}
% ==========================================================

\begin{center}
\renewcommand{\arraystretch}{1.3}

\begin{tabular}{lcc}
\toprule
\textbf{Regresión} & \textbf{Coeficiente estimado} & \textbf{Interpretación} \\
\midrule
$Y$ sobre $X^{*}$ & $\approx 2.00$ & Recupera el verdadero parámetro \\
$Y$ sobre $X$     & $\approx 1.00$ & Sesgo de atenuación \\
\bottomrule
\end{tabular}
\end{center}

Con:

\begin{align}
\Var(X^{*}) &\approx 3.92 \\
\Var(v) &\approx 4.01
\end{align}

la teoría predice:

\begin{equation}
\plim(\hat{\beta})
\approx
0.99
\end{equation}

valor extremadamente cercano al coeficiente obtenido en la simulación.

% ==========================================================
\section{Interpretación de resultados}
% ==========================================================

La simulación confirma exactamente la teoría de errores en variables.

\begin{enumerate}[leftmargin=1.8em]

\item Cuando utilizamos la variable verdadera $X_i^{*}$,
el estimador recupera correctamente:

\begin{equation}
\hat{\beta}\approx2
\end{equation}

\item Cuando utilizamos la variable observada con error:

\begin{equation}
X_i=X_i^{*}+v_i
\end{equation}

el coeficiente cae aproximadamente a:

\begin{equation}
\hat{\beta}\approx1
\end{equation}

\item El coeficiente estimado se desplaza hacia cero,
mostrando el clásico sesgo de atenuación.

\item Aunque la muestra es muy grande ($n=5000$),
el sesgo no desaparece.

\end{enumerate}

% ==========================================================
\section{Conclusión}
% ==========================================================

\begin{clave}

El error de medición en la variable explicativa provoca:

\begin{enumerate}[leftmargin=1.8em]

\item Violación del supuesto clásico:

\begin{equation*}
\Cov(X_i,u_i)=0
\end{equation*}

\item Sesgo sistemático en MCO.

\item Inconsistencia del estimador.

\item Sesgo hacia cero del coeficiente estimado
(\textit{attenuation bias}).

\end{enumerate}

El estimador no converge al verdadero parámetro poblacional:

\begin{equation*}
\plim(\hat{\beta})
=
\frac{\beta}
{1+\left(\frac{\sigma_v^2}{\sigma_X^2}\right)}
\end{equation*}

sino a una versión atenuada determinada por la razón entre:

\begin{itemize}
\item la varianza de la variable verdadera
\item la varianza del error de medición
\end{itemize}

\end{clave}

\vspace{12pt}

\noindent\rule{\linewidth}{0.4pt}

\vspace{10pt}

\begin{center}
\small\textsc{Referencias}
\end{center}

\vspace{6pt}

\noindent
Gujarati, D.~N., \& Porter, D.~C. (2010).
\textit{Econometría} (5.a ed.).
McGraw-Hill.

\vspace{6pt}

\noindent
Wooldridge, J.~M. (2016).
\textit{Introductory Econometrics} (6.a ed.).
Cengage.

\vspace{6pt}

\noindent
Greene, W.~H. (2018).
\textit{Econometric Analysis} (8.a ed.).
Pearson.

% ==========================================================
\end{document}
